\begin{statement}{C.2.UN.79}{Proved}{The \((-1,-2)\) curve with canonical orders \((0,0,1)\) occurs in a compact zero-Euler torus fibre}
There is a smooth compact projective threefold \(X\), a proper map
\(f:X\to\PP^1\) with connected general fibre a complex two-torus,
and a full reduced SNC fibre \(S=f^{-1}(0)\) with \(e(S)=0\),
containing a smooth rational curve \(F\) with the following exact
properties. The curve belongs to just one component \(P\) generically,
is not a double curve, has \(F_P^2=-1\), and meets two other
components transversely, once each, at distinct points. Its entire
analytic neighbourhood, including \(f\), is the germ along the zero
section of
\[
 \operatorname{Tot}\bigl(\OO_{\PP^1}(-1)_a
                 \oplus\OO_{\PP^1}(-2)_b\bigr),\qquad
 f=bUV,\quad [U:V]\in\PP^1.                         \tag{C.BC.1}
\]
Here \(P=(b=0)\), and the other two components are the fibres
\(U=0\) and \(V=0\) of the displayed vector bundle.
The actual generator of \(f_*K_X\) has orders \(0,1,0\)
on these three components in that order. Thus, after ordering the
two neighbours by canonical height, the tuple is \((0,0,1)\).
In particular
\[
                    N_{F/X}=\OO(-1)\oplus\OO(-2).            \tag{C.BC.2}
\]

The construction retains an actual point-blowdown map from \(P\)
to \(\PP^1\times E\), where \(E\) is an elliptic curve.
It maps \(F\) isomorphically onto \(\PP^1\times\{e_0\}\).
For the elliptic map obtained by composing this blowdown with
\(\PP^1\times E\to\PP^1\), the curve \(F\) is therefore
horizontal. The specified exceptional locus of this surface map
does not contain \(F\). This records the scope relevant to UN70:
the example concerns an arbitrary nonboundary \((-1)\)-curve,
with its exact comparison to the chosen elliptic minimalization.
\end{statement}

\begin{proof}
\par\medskip\noindent\textbf{An actual compact elliptic surface with a toric \(I_3\) chart.}
Take a compact rational elliptic surface with a simple nodal fibre.
Such a surface, with an explicit cubic-pencil construction, was
used in UN41. Alternatively the Weierstrass equation
\[
       y^2=x^3+(t^4+1)x+t^6                         \tag{C.BC.3}
\]
gives one: its discriminant, up to a nonzero constant, is
\(31t^{12}+12t^8+12t^4+4\), which is square-free and has no
zero at infinity. Choose one nodal value, put it at \(t=0\)
by a base coordinate, and take the global cubic base change
\(t=s^3\). Normalize and minimally resolve. The resulting
compact smooth projective elliptic surface \(J\to\PP^1_s\)
has an \(I_3\) fibre
\[
                         C_0+C_1+C_2 .             \tag{C.BC.4}
\]
The components are smooth rational curves, meeting in a triangle
at three distinct nodes, with no triple point. The base change
is unramified at the other singular values; at a smooth value
it is the ordinary pullback of a smooth family. Thus these
operations give a smooth total surface after resolving the
single local \(A_2\) equation at the chosen nodal value.

Here is the required local comparison, with its entire two-chart
neighbourhood retained. At the original nodal point the map is
\(xy=t\), by the holomorphic Morse coordinates. The root is
\(xy=s^3\). Two consecutive charts of its actual minimal
resolution are
\[
 \begin{aligned}
 (x,y,s)&=(b,b^2z^3,bz),\qquad (b,z)\in\C^2,\\
 (x,y,s)&=(u^2c,uc^2,uc),\qquad (u,c)\in\C^2,\\
                 u&=z^{-1},\qquad c=z^2b .
 \end{aligned}                                             \tag{C.BC.5}
\]
The curve \(C_1\) is \(b=0\) in the first chart and
\(c=0\) in the second. Its adjacent curves are \(z=0\)
and \(u=0\). The two-chart union is the actual total space
of \(\OO(-2)\), restricted to a neighbourhood if necessary;
its map is exactly \(s=bz=cu\).
The displayed monomials satisfy \(xy=s^3\), glue on their
overlap, and are the successive point-resolution charts,
so this comparison does not identify a neighbourhood merely
from its normal degree.

The two-form \(dx\wedge ds/x\) is \(db\wedge dz\)
on the first chart and \(du\wedge dc\) on the second.
An original canonical generator is a nonvanishing
holomorphic multiple of \(dx\wedge dy\) at the node.
Pull it back and divide by \(s^2\); the resulting
form is a nonvanishing holomorphic multiple of the
displayed two-form and extends through the unchanged
smooth charts of the original fibre. The three
components have the same original root order two.
This gives a nowhere-zero canonical generator
near (C.BC.4). Hence the actual
canonical evaluation orders on all three components of
\(J\) near this fibre are zero.

\par\medskip\noindent\textbf{One global curve blowup produces the required germ and orders.}
Fix an elliptic curve \(E\) and a point \(e_0\in E\).
Set \(X_0=J\times E\), with map \(f_0\) induced by \(s\).
Its general fibre is the product of two elliptic curves.
The central components are \(P_i=C_i\times E\), all
of canonical order zero. Blow up the complete smooth
rational curve
\[
                      B_0=C_0\times\{e_0\}\subset P_0 .
                                                               \tag{C.BC.6}
\]
The centre is contained generically in just \(P_0\) and
meets each of its two double curves transversely at one
point. The inverse chart of UN74, or the following direct
calculation, shows that the new fibre is full reduced SNC.
Near any such point write the original map as \(xy\),
the centre as \((x,a)\), and the blowup charts as
\(a=xA\) and \(x=aX\). Their maps are respectively
\(xy\) and \(aXy\). All displayed factors have
multiplicity one. The new exceptional component \(R_0\)
has canonical order one, because the Jacobian of this
smooth curve blowup has a simple exceptional zero.

Choose the labelling in (C.BC.5) so that \(z=0\) is
the branch of \(C_0\). Let \(a\) be a local coordinate
of \(E\) at \(e_0\). Near the whole curve
\(C_1\times\{e_0\}\), the centre is \((z,a)\)
near its first endpoint and is absent elsewhere.
In the \(z\)-chart of its blowup put
\[
               a=zA,\qquad s=zb .                          \tag{C.BC.7}
\]
The strict transform \(F\) is \(A=b=0\). In the other
chart of (C.BC.5) no modification occurs near this curve;
write its elliptic coordinate as \(a_\infty\). The actual
transition is
\[
       u=z^{-1},\qquad A=u a_\infty,\qquad b=u^2c,
             \qquad s=zb=uc .                              \tag{C.BC.8}
\]
These are exactly the two vector-bundle charts in (C.BC.1).
They prove (C.BC.2) and \(F_P^2=-1\), with the entire
map retained. The two boundary intersections occur at
\(z=0\) and \(u=0\), whereas \(F\) is contained
generically only in \(b=0\).

Up to its retained nonvanishing holomorphic factor, the actual
product canonical form pulls back as
\[
 db\wedge dz\wedge da
       =z\,db\wedge dz\wedge dA .                          \tag{C.BC.9}
\]
At the other endpoint this is a nonzero unit times
\(du\wedge dc\wedge da_\infty\). Thus it has a simple
zero exactly on \(R_0=(z=0)\), with zero orders on
\(P\) and its other old neighbour. This is the requested
canonical tuple for the actual direct-image generator.
It remains a generator: the curve blowup has the actual
canonical pushforward \(\pi_*K=K\) through its Jacobian,
or directly through Hartogs across its smooth codimension-two
centre.

The surface \(P\) is \(\operatorname{Bl}_p(\PP^1\times E)\),
where \(p\) is the endpoint of
\(\PP^1\times\{e_0\}\). The strict transform of that
rational curve is \(F\). The old double curve through
\(p\) has the other tangent direction in \(P_i\), so
its strict transform meets the exceptional line at a
different point from \(F\). In particular the full
strict transform of \(P_0\) is disjoint from \(F\).
The actual surface blowdown maps \(F\) isomorphically
to \(\PP^1\times\{e_0\}\); its composition with
the first projection has degree one on \(F\).

\par\medskip\noindent\textbf{Exact zero-Euler balancing away from the entire curve.}
We make the remaining centres explicit. Use
\[
 E:\ y^2=x(x-1)(x-2),\qquad e_0=(0,0),\quad e_1=(1,0).
                                                               \tag{C.BC.10}
\]
The strict \(P_0\) is still the identical surface
\(\PP^1\times E\): restriction of the first blowup
to this component is the blowup of its Cartier divisor
\(B_0\), hence an isomorphism. Blow up the disjoint
curve \(B_1=C_0\times\{e_1\}\). Again \(P_0\)
is unchanged and all full multiplicities are one.
The second centre is disjoint from \(F\), so this
operation preserves a neighbourhood of the whole
compact curve.

Choose the coordinate \(z\) on \(C_0\simeq\PP^1\)
so its two old boundary points are \(0,\infty\). Put
\[
 P(z)=z^2+1,\qquad Q(z)=3z^2+z+5,\qquad
 D=\{P(z)x=Q(z)\}\subset\PP^1\times E .
                                                               \tag{C.BC.11}
\]
This is the divisor of a section of
\(\OO_{\PP^1}(2)\boxtimes\OO_E(2\cdot\infty_E)\).
Its map to \(\PP^1_z\) is the exact smooth double
cover
\[
          Y^2=P Q(Q-P)(Q-2P),\qquad Y=yP^2 .                  \tag{C.BC.12}
\]
The four quadratic factors have respectively discriminants
\(-4,-59,-31,-11\), are pairwise coprime, and have
nonzero leading coefficients. Their eight simple roots
give a smooth connected genus-three curve. The chart
at \(P=0\) is the ordinary elliptic infinity chart;
since these zeros are simple it gives the same smooth
double cover, not additional components. At \(z=0\)
and \(\infty\) the values \(x=5,3\) are outside
\(\{0,1,2,\infty\}\). Thus \(D\) meets each old
elliptic boundary transversely in two points and avoids
all four existing triple points, which have elliptic
coordinate \(e_0\) or \(e_1\). The simple zeros of
\(Q\) and \(Q-P\) prove transverse intersection
with \(B_0\) and \(B_1\), respectively. These facts
also verify all its boundary intersections directly.

Blow up this \(D\) in the threefold. The same two
charts \(xy\), \(aXy\) prove reduced SNC at its
boundary crossings; elsewhere it is a smooth-centre
blowup in one reduced component. It has canonical
order one generically, since its ambient centre lies
in the zero-order component \(P_0\). The centre is
disjoint from \(F\), so (C.BC.8)--(C.BC.9) remain
unchanged on an open neighbourhood of the entire
curve, including every analytic transverse jet.

Initially \(e((C_0+C_1+C_2)\times E)=0\).
For a smooth curve centre contained in the fibre,
the actual blowup replaces that curve by its
\(\PP^1\)-bundle exceptional divisor and is an
isomorphism on the complement. Its reduced-fibre
Euler defect is therefore exactly \(e\) of the
centre. Consequently our three operations give
\[
             e(S)=0+e(B_0)+e(B_1)+e(D)
                         =0+2+2-4=0 .                       \tag{C.BC.13}
\]
For an independent stratum check, the six final components
have Euler numbers \(0,4,4,6,6,-8\). The double curves
have total Euler characteristic \(24\), and there are
\(12\) triple points. Thus their actual SNC
inclusion--exclusion gives \(12-24+12=0\), in
agreement with (C.BC.13).

All blowups are along closed smooth curves in a
compact projective threefold. Their total spaces
are smooth, compact and projective. They preserve
the map outside the single fibre; properness and
connectedness follow from their connected blowup
fibres and the original connected elliptic product.
The other fibres remain unchanged products with
\(E\). Finally, every additional point blowup
induced on \(P\) is away from \(F\), so the
same composite surface blowdown and its degree-one
map on \(F\) prove the stated horizontal
qualification on the final zero-Euler model.
\end{proof}


% C.2.UN.28 from part_C2_compact_core.tex

